% GENERATED FILE. Edit canonical lessons, then run npm run generate:books.
% canonical-source-hash: fnv1a64:7c1cf410ae0517c3
% canonical-lessons: FR-C343-avantage, FR-C343-inconvenient, FR-C343-aspect, FR-C343-critique, FR-C343-souligner

\chapter{Le pour et le contre}
\label{ch:fr-le-pour-et-le-contre}
\hlchaptermodality{343}

\begin{chapteropening}
\textbf{By the end of this chapter:} \emph{I can weigh up pros and cons: name an advantage, a drawback and an aspect, voice a criticism and stress a point.}

Complete the last lesson of chapter 343: I can weigh up pros and cons: name an advantage, a drawback and an aspect, voice a criticism and stress a point.
\end{chapteropening}

\section[l'avantage]{l'avantage — the advantage}

\label{lesson:FR-C343-avantage}



\begin{quote}
Before the new one: say the French for subjective, then the French for impartial.
\end{quote}

\subsection*{You'll want to know: l'avantage}
\textbf{l'avantage} — \textquotedblleft{}the advantage\textquotedblright{}.

\begin{quote}
Le grand avantage de ce quartier, c'est le calme.
\end{quote}

\begin{quote}
\emph{The big advantage of this neighbourhood is the quiet.}
\end{quote}

Masculine: \emph{un avantage}.

\begin{cousinweb}
\emph{avantage} is built from \emph{avant} \textquotedblleft{}before, ahead\textquotedblright{} and the suffix \emph{-age}: being ahead. \emph{avant} is usually traced to Late Latin \emph{abante} \textquotedblleft{}in front\textquotedblright{}, from \emph{ab} and \emph{ante} \textquotedblleft{}before\textquotedblright{} (Lewis \& Short).
\end{cousinweb}

\begin{grammarlens}[title={What you've built}]
One more word for telling, explaining and describing at B1.
\end{grammarlens}

\subsection*{Your turn}
\begin{itemize}
\raggedright
  \item \emph{Say it:} \emph{l'avantage}
  \item \emph{Say it:} \emph{l'avantage}, once more
  \item \emph{Say it:} a sentence of your own with \emph{l'avantage}
  \item \emph{Recall:} say the French for subjective, then the French for impartial, then say \emph{l'avantage} again
\end{itemize}

\begin{tcolorbox}[breakable,colback=teal!4,colframe=teal!35!black,title={Before you move on}]
What does l'avantage mean? (\textquotedblleft{}The advantage\textquotedblright{}.) Say it once more.
\end{tcolorbox}
\section[l'inconvénient]{l'inconvénient — the drawback, the disadvantage}

\label{lesson:FR-C343-inconvenient}



\begin{quote}
Before the new one: say the French for impartial, then the French for the advantage.
\end{quote}

\subsection*{You'll want to know: l'inconvénient}
\textbf{l'inconvénient} — \textquotedblleft{}the drawback, the disadvantage\textquotedblright{}.

\begin{quote}
Le seul inconvénient, c'est le prix.
\end{quote}

\begin{quote}
\emph{The only drawback is the price.}
\end{quote}

\emph{Je n'y vois pas d'inconvénient} — I have no objection. Masculine.

\begin{cousinweb}
\emph{inconvénient} comes from Latin \emph{inconveniēns} \textquotedblleft{}unsuitable\textquotedblright{} (Lewis \& Short), built from \emph{in-} \textquotedblleft{}not\textquotedblright{} and \emph{convenīre} \textquotedblleft{}to come together, to suit\textquotedblright{}, on \emph{venīre} \textquotedblleft{}to come\textquotedblright{}.
\end{cousinweb}

\begin{grammarlens}[title={What you've built}]
One more word for telling, explaining and describing at B1.
\end{grammarlens}

\subsection*{Your turn}
\begin{itemize}
\raggedright
  \item \emph{Say it:} \emph{l'inconvénient}
  \item \emph{Say it:} \emph{l'inconvénient}, once more
  \item \emph{Say it:} a sentence of your own with \emph{l'inconvénient}
  \item \emph{Recall:} say the French for impartial, then the French for the advantage, then say \emph{l'inconvénient} again
\end{itemize}

\begin{tcolorbox}[breakable,colback=teal!4,colframe=teal!35!black,title={Before you move on}]
What does l'inconvénient mean? (\textquotedblleft{}The drawback, the disadvantage\textquotedblright{}.) Say it once more.
\end{tcolorbox}
\section[l'aspect]{l'aspect — the aspect, the side (of a question)}

\label{lesson:FR-C343-aspect}



\begin{quote}
Before the new one: say the French for the advantage, then the French for the drawback.
\end{quote}

\subsection*{You'll want to know: l'aspect}
\textbf{l'aspect} — \textquotedblleft{}the aspect, the side (of a question)\textquotedblright{}.

\begin{quote}
Il faut regarder tous les aspects du problème.
\end{quote}

\begin{quote}
\emph{We have to consider every aspect of the problem.}
\end{quote}

The final \emph{ct} is usually silent. Masculine.

\begin{cousinweb}
\emph{aspect} is a learned borrowing from Latin \emph{aspectus} \textquotedblleft{}a seeing, a look, an appearance\textquotedblright{} (EtymDB; Lewis \& Short), from \emph{aspicere} \textquotedblleft{}to behold\textquotedblright{}, built on \emph{specere} \textquotedblleft{}to look\textquotedblright{}.
\end{cousinweb}

\begin{grammarlens}[title={What you've built}]
One more word for telling, explaining and describing at B1.
\end{grammarlens}

\subsection*{Your turn}
\begin{itemize}
\raggedright
  \item \emph{Say it:} \emph{l'aspect}
  \item \emph{Say it:} \emph{l'aspect}, once more
  \item \emph{Say it:} a sentence of your own with \emph{l'aspect}
  \item \emph{Recall:} say the French for the advantage, then the French for the drawback, then say \emph{l'aspect} again
\end{itemize}

\begin{tcolorbox}[breakable,colback=teal!4,colframe=teal!35!black,title={Before you move on}]
What does l'aspect mean? (\textquotedblleft{}The aspect, the side (of a question)\textquotedblright{}.) Say it once more.
\end{tcolorbox}
\section[la critique]{la critique — the criticism; the review}

\label{lesson:FR-C343-critique}



\begin{quote}
Before the new one: say the French for the drawback, then the French for the aspect.
\end{quote}

\subsection*{You'll want to know: la critique}
\textbf{la critique} — \textquotedblleft{}the criticism; the review\textquotedblright{}.

\begin{quote}
Elle n'aime pas les critiques.
\end{quote}

\begin{quote}
\emph{She does not like criticism.}
\end{quote}

\emph{Une critique} is a criticism or a review; \emph{un critique} is a critic. The verb is \emph{critiquer}.

\begin{cousinweb}
\emph{critique} comes from Latin \emph{criticus}, which came from Greek \emph{kritikós} \textquotedblleft{}able to judge\textquotedblright{} (EtymDB), built on \emph{krínein} \textquotedblleft{}to separate, to judge\textquotedblright{}.
\end{cousinweb}

\begin{grammarlens}[title={What you've built}]
One more word for telling, explaining and describing at B1.
\end{grammarlens}

\subsection*{Your turn}
\begin{itemize}
\raggedright
  \item \emph{Say it:} \emph{la critique}
  \item \emph{Say it:} \emph{la critique}, once more
  \item \emph{Say it:} a sentence of your own with \emph{la critique}
  \item \emph{Recall:} say the French for the drawback, then the French for the aspect, then say \emph{la critique} again
\end{itemize}

\begin{tcolorbox}[breakable,colback=teal!4,colframe=teal!35!black,title={Before you move on}]
What does la critique mean? (\textquotedblleft{}The criticism; the review\textquotedblright{}.) Say it once more.
\end{tcolorbox}
\section[souligner]{souligner — to stress, to point out; to underline}

\label{lesson:FR-C343-souligner}



\begin{quote}
Before the new one: say the French for the aspect, then the French for the criticism.
\end{quote}

\subsection*{You'll want to know: souligner}
\textbf{souligner} — \textquotedblleft{}to stress, to point out; to underline\textquotedblright{}.

\begin{quote}
Il a souligné deux erreurs dans mon texte.
\end{quote}

\begin{quote}
\emph{He pointed out two mistakes in my text.}
\end{quote}

Literally to draw a line under: \emph{souligner un mot}.

\begin{cousinweb}
\emph{souligner} is built from \emph{sous} \textquotedblleft{}under\textquotedblright{} and \emph{ligne} \textquotedblleft{}line\textquotedblright{}, which comes from Latin \emph{līnea} \textquotedblleft{}a linen thread, a line\textquotedblright{} (Lewis \& Short).
\end{cousinweb}

\begin{grammarlens}[title={What you've built}]
One more word for telling, explaining and describing at B1.
\end{grammarlens}

\subsection*{Your turn}
\begin{itemize}
\raggedright
  \item \emph{Say it:} \emph{souligner}
  \item \emph{Say it:} \emph{souligner}, once more
  \item \emph{Say it:} a sentence of your own with \emph{souligner}
  \item \emph{Recall:} say the French for the aspect, then the French for the criticism, then say \emph{souligner} again
\end{itemize}

\begin{tcolorbox}[breakable,colback=teal!4,colframe=teal!35!black,title={Before you move on}]
What does souligner mean? (\textquotedblleft{}To stress, to point out; to underline\textquotedblright{}.) Say it once more.
\end{tcolorbox}
